% ===================================================================
%  ТАБЛИЦЯ № 14 — Вулиця. — Торговці.
%  Traduction 2026 d'apres texto/io, controlee sur texto/fr, texto/en
%  et texto/es. Consigne : voir texto/uk/10-tablycia-01.tex.
%
%  LE TABLEAU QUI AVAIT COUTE DEUX INVERSIONS AU SUEDOIS. Le jet (49)
%  de la fontaine (48) au bloc 14-1, l'imperiale (59) de l'omnibus (58)
%  au bloc 15-2 : dans les deux cas le possede porte le numero, et une
%  langue qui prepose son genitif renverse les renvois. Le suedois
%  avait du echanger le genitif contre une preposition. L'ukrainien
%  postpose le sien comme l'ido — « струмінь (49) фонтана (48) »,
%  « імперіалі (59) омнібуса (58) » — et ne paie rien.
%
%  DEUX HOMONYMES A EVITER. La dactylo du bloc 11-1 est un
%  « машиніст » en russe, mais ce mot a deja servi au tableau 12 pour
%  le mecanicien de la locomotive, et objekti.py releve les objets par
%  leur nom. On ecrit donc « переписувач ».
% ===================================================================

%%K t14-apar-1 apar
\VUsaut{4.0mm}
\VUtitre{15.0pt}{60}{ТАБЛИЦЯ № 14}
\VUsaut{3.0mm}

%%K t14-tit-1 sub
\VUpk{11.4pt}{40}{Вулиця. --- Торговці.}
\VUsaut{3.0mm}

%%K t14-01-1 p
1. --- Мій друг Іван, який живе на селі, приїхав провести три дні в моїх.
\VUgras{До} цього часу йому жодного разу не траплялося бачити великого міста,
тож ми цілими днями гуляємо, і я пояснюю йому все, чого він не знає.

%%K t14-02-1 p
2. --- Спершу ми пішли дивитися \VUgras{обсерваторію} \textsuperscript{(1)}, яка
дуже його зайняла. Вертаючись \VUgras{вулицею} \textsuperscript{(3)}, де
містяться \VUgras{турецькі лазні} \textsuperscript{(2)}, ми стріли
\VUgras{похорон} \textsuperscript{(4)}. На чолі \VUgras{похоронної ходи} йшло
\VUgras{духівництво}, поперед \VUgras{катафалка}, на який поставили
\VUgras{труну}. Багато друзів ішло за родиною в \VUgras{жалобі}.

%%K t14-02-2 p
Коли хода пройшла, ми перейшли вулицю перед великою \VUgras{пивницею}
\textsuperscript{(5)}, де багато \VUgras{відвідувачів} пило \VUgras{пиво} на
терасі під заскленим \VUgras{дашком} \textsuperscript{(6)}.

%%K t14-03-1 p
3. --- Івана спантеличив \VUgras{герб}, поміщений між крамницею
\VUgras{торговця винами й горілками} \textsuperscript{(7)} і крамницею
\VUgras{продавця нот} \textsuperscript{(10)}. Він спитав мене, що це значить; я
сказав йому, що це англійське \VUgras{консульство} \textsuperscript{(8)}, і
показав \VUgras{консула} \textsuperscript{(9)}, що стояв на балконі.

%%K t14-tit-2 sub
\VUpk{10.2pt}{40}{Розглядаючи вітрини.}
\VUsaut{3.0mm}

%%K t14-04-1 p
4. --- Певна річ, ми спинилися перед виставкою \VUgras{музичних інструментів},
\VUgras{фісгармоній}, \VUgras{органів} і \VUgras{роялів}; ми розглядали
ілюстровані обкладинки \VUgras{нот} і фортепіанних \VUgras{п'єс} і милувалися
золоченою дерев'яною \VUgras{лірою}, що правила за вивіску.

%%K t14-05-1 p
5. --- Хмари пилу, зняті \VUgras{двірниками} \textsuperscript{(11)} і
\VUgras{замітальною машиною} \textsuperscript{(12)}, змусили нас швидше пройти
повз прекрасні \VUgras{меблеві гарнітури} в \VUgras{мебляра}
\textsuperscript{(13)} і повз виставку \VUgras{печей} і \VUgras{ламп} у
\VUgras{залізній крамниці} \textsuperscript{(14)}.

%%K t14-06-1 p
6. --- У цьому самому місці нам і справді довелося зійти з тротуару, бо він був
захаращений паками, які знімали з \VUgras{меблевого фургона}
\textsuperscript{(16)}, щоб унести до щойно винайнятої \VUgras{крамниці}
\textsuperscript{(15)}.

%%K t14-06-2 p
Це завадило нам побачити прекрасні екіпажі в \VUgras{каретника}
\textsuperscript{(17)}, але не завадило спинитися й подивитися, як
\VUgras{пакувальник} \textsuperscript{(19)} збиває ящик для свого сусіда,
\VUgras{торговця велосипедами й автомобілями} \textsuperscript{(18)}. Потім, що
було вже доволі пізно, ми вернулися додому.

%%K t14-tit-3 sub
\VUpk{10.2pt}{40}{Прогулянка містом.}
\VUsaut{3.0mm}

%%K t14-07-1 p
7. --- Другого дня мати запропонувала, щоб Іван знявся, і попросила мене
провести його до \VUgras{фотографа} \textsuperscript{(20)}. Після сніданку ми
пішли до \VUgras{ательє} художника, де довелося чимало почекати. Щоб згаяти час,
ми розглядали \VUgras{портрети} \textsuperscript{(22)}, що висіли на стіні.

%%K t14-07-2 p
Коли прийшла наша черга, справа вийшла недовга, і, ледве мій друг скінчив
позувати перед \VUgras{апаратом} \textsuperscript{(21)}, ми спустилися.

%%K t14-08-1 p
8. --- На першому поверсі ми побачили крізь відчинені двері \VUgras{перукарні}
\textsuperscript{(23)}, як підмайстер стриг клієнта.

%%K t14-08-2 p
Знову вийшовши на вулицю, ми пройшли повз \VUgras{тютюнову крамницю}
\textsuperscript{(24)}. Івана так потішили червоний \VUgras{ліхтар}
\textsuperscript{(25)} і \VUgras{велика люлька}, що правила за \VUgras{вивіску}
\textsuperscript{(26)}, що він налетів на двох стариків, які розмовляли за
\VUgras{понюшкою тютюну} \textsuperscript{(27)}.

%%K t14-tit-4 sub
\VUpk{10.2pt}{40}{Цукерня.}
\VUsaut{3.0mm}

%%K t14-09-1 p
9. --- \VUgras{Асигнації} і \VUgras{монети} всіх країн, виставлені у
\VUgras{вітрині} \VUgras{міняйла} \textsuperscript{(29)}, затримали нас на
хвилину, але що була година підвечірка, ми, не гаючись довше, увійшли до
\VUgras{цукерні} \textsuperscript{(31)}.

%%K t14-10-1 p
10. --- Побачивши всі \VUgras{ласощі}, які були в крамниці --- \VUgras{тістечка,
медяники, печиво} і \VUgras{цукерки} всякого роду --- ми насилу могли вибрати.
Кінець кінцем Іван спинився на \VUgras{кремових тістечках}, а я взяв тільки
маленький \VUgras{вишневий торт}, бо від солодкого \VUgras{у мене болять зуби},
а мені зовсім не хочеться надто часто ходити до \VUgras{зубного лікаря}
\textsuperscript{(30)}.

%%K t14-tit-5 sub
\VUpk{10.2pt}{40}{Писання на машинці.}
\VUsaut{3.0mm}

%%K t14-11-1 p
11. --- Щоб показати другові \VUgras{друкарську машинку}, про яку він чув, але
якої не бачив, ми зайшли до \VUgras{контори} \textsuperscript{(32)} мого
\VUgras{двоюрідного брата} \textsuperscript{(34)}. Ми застали його, як він
диктував листи своїй \VUgras{секретарці} \textsuperscript{(35)}, яка
\VUgras{стенографістка}. Ледве лист був скінчений, \VUgras{переписувач}
\textsuperscript{(33)} переписував його на своїй машинці. Не хотівши заважати
роботі родича, ми пробули в нього лише кілька хвилин.

%%K t14-12-1 p
12. --- Під конторою двоюрідного брата живе великий \VUgras{годинникар і ювелір}
\textsuperscript{(37)}, який до того ж \VUgras{срібляр}. У його чудовій крамниці
багато \VUgras{стінних годинників, камінних годинників, кишенькових
годинників}, \VUgras{ланцюжків} і дорогих \VUgras{коштовностей}: \VUgras{перлових
намист}, золотих \VUgras{браслетів}, \VUgras{сережок} і \VUgras{перснів} із
\VUgras{самоцвітами} і \VUgras{діамантами}. Одну з вітрин віддано цілком
\VUgras{сріблу}, і вона вміщає велику збірку срібних \VUgras{приборів}.

%%K t14-13-1 p
13. --- Завертаючи за ріг вулиці, я помітив, що \VUgras{табличку}
\textsuperscript{(36)} поруч із \VUgras{номером} дому перемінили. Я збирався
сказати це вголос, коли залунали веселі звуки військового оркестру. Ми
пустилися бігти до полку, не подумавши, що він пройде повз крамниці
\VUgras{капелюшника} \textsuperscript{(39)} і \VUgras{рукавичника}
\textsuperscript{(40)} і що ми були б так само добре поставлені, щоб дивитися на
його проходження, якби зосталися перед вітриною \VUgras{оптика}
\textsuperscript{(38)}, повною \VUgras{пенсне, окулярів} і \VUgras{біноклів}.

%%K t14-tit-6 sub
\VUpk{10.2pt}{40}{Проходження полку.}
\VUsaut{3.0mm}

%%K t14-14-1 p
14. --- На чолі \VUgras{полку} \textsuperscript{(41)} йшов \VUgras{тамбурмажор}
\textsuperscript{(42)}, за ним \VUgras{барабанщики} \textsuperscript{(43)},
\VUgras{сурмачі} \textsuperscript{(44)} і весь \VUgras{оркестр}.
\VUgras{Генерал} \textsuperscript{(45)}, який був на площі зі своїм ад'ютантом,
спинився біля \VUgras{струменя} \textsuperscript{(49)} \VUgras{фонтана}
\textsuperscript{(48)}, щоб дивитися \VUgras{парад}.

%%K t14-14-2 p
\VUgras{Штабні офіцери} \textsuperscript{(46)}, \VUgras{полковник} і
\VUgras{підполковник} салютували йому шпагами. \VUgras{Майори} йшли на чолі
своїх \VUgras{батальйонів} \textsuperscript{(47)}, і кожен \VUgras{капітан}
командував своєю \VUgras{ротою}, а \VUgras{поручники}, \VUgras{підпоручники} і
\VUgras{унтер-офіцери} займали свої звичайні місця поруч зі своїми людьми.

%%K t14-15-1 p
15. --- Багато перехожих зібралося на \VUgras{перехресті}
\textsuperscript{(50)}; крамарі --- \VUgras{парасольник} \textsuperscript{(51)},
\VUgras{фарбар} \textsuperscript{(53)}, \VUgras{зброяр} \textsuperscript{(54)} ---
вийшли подивитися на \VUgras{солдатів}. Усі екіпажі --- \VUgras{візницькі
прольотки} \textsuperscript{(52)}, \VUgras{власні екіпажі}
\textsuperscript{(57)}, а також \VUgras{поливальна бочка} \textsuperscript{(56)}
--- притислися до тротуару.

%%K t14-tit-7 sub
\VUpk{10.2pt}{40}{Сцени на вулиці.}
\VUsaut{3.0mm}

%%K t14-15-2 p
\VUgras{Комівояжер} \textsuperscript{(55)}, що ніс у руках свої \VUgras{валізи зі
зразками}, жваво перейшов вулицю, а ті, що сиділи на \VUgras{імперіалі}
\textsuperscript{(59)} \VUgras{омнібуса} \textsuperscript{(58)}, оберталися, щоб
помилуватися видовищем. Бідний \VUgras{вуличний співак} \textsuperscript{(60)} зі
своєю \VUgras{катеринкою} \textsuperscript{(61)} мусив урвати свою
\VUgras{пісню}, бо гуркіт барабанів заглушав його голос.

%%K t14-16-1 p
16. --- Коли полк порівнявся із \VUgras{суконною крамницею}
\textsuperscript{(64)}, \VUgras{швачки} \textsuperscript{(63)} і \VUgras{кравці}
\textsuperscript{(62)} покинули свої \VUgras{швейні машини} і роботу, щоб
виглянути у вікно. \VUgras{Крамар} \textsuperscript{(65)}, який щойно виставив
нові \VUgras{костюми} \textsuperscript{(66)} у свою \VUgras{вітрину}, мусив
прибрати штуки \VUgras{сукна} \textsuperscript{(67)} і \VUgras{полотна},
розкладені на тротуарі коло \VUgras{водостоку} \textsuperscript{(68)}, боячись,
що натовп їх поперекидає; навіть старому \VUgras{рознощикові}
\textsuperscript{(69)}, у якого воротарка торгувала панчохи, довелося зайти в
підворіття, щоб докінчити продаж.

%%K t14-16-2 p
Ми провели полк до касарень \VUgras{і}, дуже вдоволені своїм днем, вернулися
додому на обід.

%%K t14-tit-8 sub
\VUpk{10.2pt}{40}{Різні ремесла й заняття.}
\VUsaut{3.0mm}

%%K t14-17-1 p
17. --- Другого дня батько запропонував звести нас до друкарні місцевої газети.
Ми пристали із захватом.

%%K t14-17-2 p
\VUgras{Друкар} --- він же й \VUgras{книгар} \textsuperscript{(70)},
\VUgras{видавець}, \VUgras{паперовик} і \VUgras{палітурник}. Він помістив на
першому поверсі \VUgras{редакцію} \textsuperscript{(71)} газети, а також
\VUgras{граверню}, де працюють \VUgras{літографи} \textsuperscript{(72)} і
\VUgras{гравери по міді} \textsuperscript{(73)}, і, нарешті, складальню, де сидять
\VUgras{складачі} \textsuperscript{(74)}. Власне \VUgras{машинна зала}
\textsuperscript{(75)}, \VUgras{друкарські верстати} і різні машини --- у нижньому
поверсі.

%%K t14-tit-9 sub
\VUpk{10.2pt}{40}{Іван ставить безліч запитань.}
\VUsaut{3.0mm}

%%K t14-18-1 p
18. --- Вийшовши з друкарні, Іван спинився у великому подиві перед скляною
скринькою на стовпі проти \VUgras{галантерейної крамниці} \textsuperscript{(77)}
і спитав, нащо вона. Батько пояснив, що це \VUgras{пожежний сповіщувач}
\textsuperscript{(76)}. І справді, все в місті було для мого друга дивом: і
простий \VUgras{питний фонтанчик} \textsuperscript{(78)}, і легкий
\VUgras{розвізний трицикл} \textsuperscript{(80)}, на якому їхав
\VUgras{посильний} \textsuperscript{(79)} у лівреї, і \VUgras{поштовий фургон}
\textsuperscript{(81)}.

%%K t14-19-1 p
19. --- Поки батько лишив нас на хвилину, щоб поговорити зі \VUgras{збирачем
податків} \textsuperscript{(83)}, який спинився порозмовляти з
\VUgras{адвокатом} \textsuperscript{(82)} і \VUgras{нотарем}
\textsuperscript{(84)}, ми бавилися, стежачи за рухом на вулиці. Повз нас
проходив найрізніший народ: \VUgras{цукерничий хлопець} \textsuperscript{(85)},
що ніс у кошику чудовий \VUgras{пиріг}, ішов за \VUgras{лахмітником}
\textsuperscript{(86)}; \VUgras{ліхтарник} \textsuperscript{(87)} розмовляв із
\VUgras{газівником} \textsuperscript{(88)}; \VUgras{черниця}
\textsuperscript{(89)}, що тримала за руку малу сирітку, верталася до свого
монастиря.

%%K t14-tit-10 sub
\VUpk{10.2pt}{40}{На зворотному шляху.}
\VUsaut{3.0mm}

%%K t14-20-1 p
20. --- Якраз коли батько до нас вернувся, Іван зареготав, побачивши замурзані
обличчя і дивне вбрання двох \VUgras{сажотрусів} \textsuperscript{(91)}, чорних
від сажі.

%%K t14-21-1 p
21. --- Я хотів купити матері рослину в \VUgras{квіткарки}
\textsuperscript{(92)}, але батько зауважив, що її важко буде нести і що краще
взяти \VUgras{помаранчів}. Ми підійшли до \VUgras{торговки з лотка}
\textsuperscript{(94)}, і, коли \VUgras{гувернантка} \textsuperscript{(93)}, яка
вибирала їх своєму вихованцеві, скінчила купувати, нас обслужили.

%%K t14-22-1 p
22. --- На зворотному шляху ми стріли кількох знайомих: стару приятельку нашої
родини, що спиралася на руку своєї \VUgras{компаньйонки}
\textsuperscript{(95)}, \VUgras{інженера} \textsuperscript{(96)} шляхів
сполучення і одну з моїх двоюрідних сестер з її \VUgras{нянею}
\textsuperscript{(98)}.

%%K t14-22-2 p
Потім ми пройшли повз \VUgras{модистку} \textsuperscript{(97)} моєї матері і
двох \VUgras{ченців} \textsuperscript{(99)}, приїжджих у місті, які підійшли
спитати в батька дорогу на вокзал.
\VUsaut{6.0mm}
\VUfilet{22mm}
